\begin{figure}[H]
\centering
\begin{tikzpicture}[font=\small,>=Latex,xscale=1.25,yscale=1.12]
  \fill[orange!8] (-1.6,-1.35) rectangle (1.6,1.35);
  \draw[gray] (-1.6,-1.35) rectangle (1.6,1.35);
  \draw[thick,blue!65!black] (-1.6,-1.35)--(1.6,-1.35);
  \draw[thick,blue!65!black] (-1.6,1.35)--(1.6,1.35);
  \draw[blue!65!black,dashed] (-1.6,-1.35)--(-.92,-.13);
  \draw[blue!65!black,dashed] (1.6,-1.35)--(.92,-.13);
  \draw[blue!65!black,dashed] (-1.6,1.35)--(-.92,.13);
  \draw[blue!65!black,dashed] (1.6,1.35)--(.92,.13);
  \fill[orange!45] (-.38,-.13) rectangle (.38,.13);
  \draw[thick] (-.38,-.13) rectangle (.38,.13);
  \draw[->,gray] (-1.8,0)--(1.88,0) node[right,black] {\(\xi\)};
  \draw[->,gray] (0,-1.70)--(0,1.75) node[above,black] {\(\theta\)};
  \node[anchor=east] at (-1.67,-1.35) {\(-r\)};
  \node[anchor=east] at (-1.67,1.35) {\(r\)};
  \node[fill=white,inner sep=1pt] at (0,1.05) {upper data};
  \node[fill=white,inner sep=1pt] at (0,-1.05) {lower data};
  \draw[->] (.9,.28)--(.34,.08);
  \node[anchor=west] at (.72,.35) {pulse};
\end{tikzpicture}
\caption{The hyperbolic slab, schematically and not to scale.
Cap regularity gives smooth data on the two cuts \(\theta=\pm r\).
The dashed lines indicate the shrinking spatial intervals used to
propagate those data toward both pulse edges. The pulse is supported
in \(\abs\theta<\delta/\tau\) and oscillates in \(\xi\).
The entire slab lies in negative curvature.}
\label{fig:pulse}
\end{figure}
